\subsection*{全体概要：この章で学ぶこと}
\addcontentsline{toc}{subsection}{全体概要：この章で学ぶこと}
この章は、ソフトウェア開発で使われる「モデル」と「方法」を扱う。モデルとは、対象を理解し、判断し、関係者へ説明するための抽象表現である。方法とは、開発を体系的に進めるための考え方、手順、成果物、確認方法のまとまりである。

SWEBOK 第11章は、次の四つの領域に分けて説明している。
\begin{enumerate}
  \item \term{モデリング}{Modeling}：モデルとは何か、どのような原則で作るか、構文・意味・実用上の読み方をどう扱うかを学ぶ。
  \item モデルの種類：構造モデルと振る舞いモデルを中心に、モデルが何を表すかを学ぶ。
  \item \term{モデルの分析}{Model Analysis}：完全性、整合性、正しさ、追跡可能性、相互作用を確認する方法を学ぶ。
  \item \term{ソフトウェア工学方法}{Software Engineering Method}：経験則に基づく方法、形式手法、プロトタイピング、アジャイル方法を学ぶ。
\end{enumerate}

\noindent 到達目標\par
この章を読み終えると、次のことができるようになる。
\begin{itemize}
  \item モデルを「対象の縮小版」ではなく「目的を持つ抽象表現」として説明できる。
  \item 構造モデルと振る舞いモデルの違いを説明できる。
  \item 構文、意味論、語用論の違いを説明できる。
  \item 事前条件、事後条件、不変条件を簡単な例で説明できる。
  \item モデルの完全性、整合性、正しさ、追跡可能性、相互作用を確認する観点を説明できる。
  \item 経験則的方法、形式手法、プロトタイピング、アジャイル方法の使い分けを説明できる。
\end{itemize}

\par\smallskip
\begin{center}
\centering
\IfFileExists{assets/generated_figures/ch11/fig-ch11-01.png}{\includegraphics[width=0.96\linewidth]{assets/generated_figures/ch11/fig-ch11-01.png}}{\fbox{\parbox[c][48mm][c]{0.9\linewidth}{\centering 画像生成待ち\\第11章の全体地図}}}
\par\smallskip
\refstepcounter{figure}\label{fig:ch11-01}図 \thefigure: 第11章の全体地図
\end{center}
図 \thefigure は、第11章の全体地図を視覚的に整理したものである。
\par\smallskip
